% ============================================================
\subsection{Objetivo 1. Diseñar un prototipo funcional de sistema de gestión nutricional que integre datos antropométricos y bioquímicos del paciente con un catálogo de alimentos ecuatorianos, empleando un modelo de lenguaje largo con base de conocimiento nutricional estructurada}
% ============================================================

Para el desarrollo del Objetivo~1 se siguieron las fases de la metodología
RAD~\cite{fauzi2023} definidas en la sección de Metodología, organizadas en
cuatro actividades según el cronograma del proyecto:

\begin{enumerate}[leftmargin=2em]
  \item \textbf{Planificación de requisitos e identificación del catálogo
        alimentario ecuatoriano} (semanas 1--3): especificación de requisitos
        funcionales y no funcionales, identificación de actores, modelo
        conceptual y selección de las fuentes del catálogo de alimentos.
  \item \textbf{Definición de la arquitectura C4 y configuración del RAG}
        (semanas 3--7): diagramas de casos de uso, modelo de dominio y de
        comportamiento, arquitectura en los cuatro niveles del modelo
        C4~\cite{brown2024c4} y configuración del pipeline de generación
        aumentada por recuperación~\cite{lewis2020}.
  \item \textbf{Desarrollo de la interfaz del sistema} (semanas 6--9):
        construcción del backend, del frontend y de la base de datos, junto con
        las pruebas de caja negra de la API.
  \item \textbf{Documentación de resultados} (semanas 9--10): consolidación de
        los artefactos producidos, documentación técnica y manual de usuario.
\end{enumerate}

Como resultado se obtuvo el prototipo web \textbf{NutriAPP}, cuyos artefactos
se describen en las secciones siguientes.
